\documentclass[10pt,a4paper]{amsart}
\usepackage[margin=19mm]{geometry}
\usepackage{amsmath,amssymb,mathtools}
\usepackage[scr=rsfs,cal=euler]{mathalfa}
\usepackage{xfrac}
\usepackage[unicode,hidelinks,bookmarks=false]{hyperref}
\newcommand{\ie}{i.e.\ }
\newcommand{\cf}{cf.\ }
\newcommand{\resp}{respectively\ }
\newcommand{\Subsection}{\subsection}
\providecommand{\nobreakdash}{\nobreak\hbox{-}}
\setlength{\emergencystretch}{2em}
\multlinegap0pt
\usepackage[shortlabels]{enumitem}
\usepackage{xcolor}
\usepackage{amssymb}
\usepackage{mathtools}
\usepackage{tikz}
\newtheorem{lemma}{Lemma}
\newtheorem{claim}{Claim}
\newtheorem{claimproof}{Claimproof}
\newtheorem{fact}{Fact}
\title{Towards the Lov\'{a}sz conjecture via sublinear expanders}
\author{Matija Buci{\'c} \and Micha Christoph \and Alexey Pokrovskiy \and Raphael Steiner}
\date{Source: arXiv:2606.09742v2 (CC BY 4.0)}
\begin{document}
\maketitle

\noindent\emph{Source and licence.} Towards the Lov\'{a}sz conjecture via sublinear expanders. Version arXiv:2606.09742v2 (CC BY 4.0), distributed by arXiv under the Creative Commons Attribution 4.0 International licence, http://creativecommons.org/licenses/by/4.0/, as recorded in the arXiv licence metadata for this exact version and shown on the abstract page https://arxiv.org/abs/2606.09742v2. Full licence notice: \texttt{licenses/arxiv-2606.09742-CC-BY-4.0.txt}.\par\smallskip

\noindent\textit{Editorial scope.} Counted unit: the article's lemma cycles (source0.tex lines 298--312). The article's complete original proof of it is delivered with it (source0.tex lines 313--394, one \texttt{proof} environment of 82 lines). No mathematics is written, repaired or re-derived: the statement and the proof are the article's own lines, in the article's own notation, and no retained line is substituted or re-flowed.  No further source-local context is needed: every cross-reference of every retained line resolves inside the two delivered spans.  Nothing was omitted from any retained span, so the package carries no omission record.  No retained line contains a citation, so the wrapper carries no bibliography and no bibliography entry.  No retained line contains a figure environment, an image-inclusion command or drawing code, so no image is delivered and the package declares no asset dependency.  Editorial formatting only, in addition: an AMS \texttt{amsart} preamble that restates the article's own declarations and macros used by the retained lines, namely the environments \texttt{lemma}, \texttt{claim}, \texttt{claimproof}, \texttt{fact} on separate counters, together with the standard packages those lines need.  The retained environments print in this document as Lemma 1, Claim 1, Claimproof 1, Fact 1 in order of appearance, whereas the article prints the same items with the numbering of its own full document.  The complete native source of the article as distributed in the arXiv e-print for this exact version is bundled unchanged as \texttt{sources/202424/source0.tex}.
\begin{lemma}\label{lem: root cycles}
    Let $\mathcal C$ be a collection of edge-disjoint cycles in some graph, let $w:V(\mathcal{C})\rightarrow \mathbb{Z}_{\ge 0}$ be a non-negative integer weight function on its vertices, and let $\mathcal R \subseteq \mathcal C$ be a subcollection of cycles such that $\mathcal R$ contains at least one cycle in every connected component of $\mathcal C$. Let $k\in \mathbb{N}$, $\delta \in (0,1)$ and $p\in [k^{-\delta},1]$. Then, there exists another subcollection $\mathcal H\subseteq \mathcal C$ meeting each of the following conditions:
    \begin{enumerate}[label=(C\arabic*), ref=(C\arabic*)]
    \item \label{cond:C1} $|\mathcal{H}\cap \mathcal{R}|\le k$,
    \item \label{cond:C2} each cycle in $\mathcal H\cap \mathcal R$ intersects at most $2k$ cycles in $\mathcal H$ in some vertex,
    \item \label{cond:C3} every cycle in $\mathcal H$ intersects at most $3k$ cycles in $\mathcal H$ in some vertex, and, 
    \item \label{cond:C4}\textcolor{white}{hack}

    \vspace{-1.4cm}
    \begin{align}\label{exp: cycle expectation}
        \mathbb E[w(V)]\geq w(\mathcal C)^{1-\delta},
    \end{align}
    where $V$ is the set of vertices contained in connected components of $\mathcal C_p\cup \mathcal H$ which contain at least one cycle of $\mathcal R$.
\end{enumerate}
\end{lemma}

\begin{proof}
     We proceed by induction on $|\mathcal C|$, where the base case ($|\mathcal C|=0$) is trivial. Now suppose we are given a non-empty collection $\mathcal{C}$ of edge-disjoint cycles in some graph and that we have proved the assertion of the lemma for all edge-disjoint cycle collections smaller than $\mathcal{C}$. Let $w,\mathcal{R},k,\delta,p$ be given as in the lemma statement. Our goal in the following is to carry out the induction step, i.e.\ to show that the assertion of the lemma holds for $\mathcal{C}$.

     Let $t=|\mathcal R|$. Let $\mathcal K_1,\ldots,\mathcal K_t$ be disjoint subsets of $\mathcal C$ obtained by iteratively selecting $\mathcal K_i\subseteq \mathcal C\setminus(\cup_{j<i}\mathcal K_j)$ to be a set of cycles lexicographically maximizing $(w(V(\mathcal K_i)\setminus \bigcup_{j<i}V(\mathcal{K}_j)),|\mathcal K_i|)$ under the constraints that $\mathcal K_i$ is connected and $\mathcal K_i$ contains exactly one cycle of $\mathcal R$. Let $R_1,\ldots, R_t$ be the cycles in $\mathcal R$ such that $R_i\in \mathcal K_i$. The following claim introduces a central observation about this partition, which will be used repeatedly.
     \begin{claim}\label{clm: disjointness}
     We have that $\mathcal{K}_1,\ldots,\mathcal{K}_t$ form a partition of $\mathcal C$. Furthermore,  $V(\mathcal K_i)\cap V(\mathcal K_j\setminus \{R_j\})=\emptyset$ for all $1\le i<j\le t$.
     \end{claim}
     \begin{claimproof}
For the first part, suppose towards a contradiction that $\mathcal C\setminus \bigcup_{i=1}^t \mathcal K_i\neq \emptyset$. Since by assumption in the lemma, every component of $\mathcal C$ contains a cycle in $\mathcal R$, and since furthermore every cycle in $\mathcal R$ belongs to exactly one of $\mathcal K_1,\ldots,\mathcal K_t$, it now follows that there must exist some $C\in \mathcal C\setminus \bigcup_{i=1}^t \mathcal K_i$ such that $V(C)\cap V(\mathcal K_i)\neq \emptyset$ for some $1\le i \le t$. Then, $\mathcal{K}_i\cup \{C\}$ is also connected, still contains precisely one cycle from $\mathcal R$, and is disjoint from $\bigcup_{k<i}\mathcal K_k$. Hence, we obtain a contradiction to the maximality of $|\mathcal K_i|$ in our choice. 

The proof of the second part of the claim is analogous:
 Suppose towards a contradiction that there exists $C\in \mathcal K_j\setminus\{R_j\}$ such that $V(C)\cap V(\mathcal K_i)\neq \emptyset$. Then, again, $\mathcal K_i\cup \{C\}$ is connected, contains precisely one cycle from $\mathcal{R}$, and is still disjoint from $\bigcup_{k<i}\mathcal{K}_k$. So, as before, we obtain a contradiction to our choice of $\mathcal{K}_i$.
     \end{claimproof}

    For all $i\in [t]$, let us set $\mathcal K_i^-:= \mathcal K_i\setminus \{R_i\}$ and let $\mathcal R_i\subseteq \mathcal K_i^-$ denote the set of all cycles in $\mathcal{K}_i^-$ which intersect $R_i$ in at least one vertex. Let further $\mathcal K_{\leq k}:=\bigcup_{i=1}^{\min\{k,t\}} \mathcal K_i$, $\mathcal K_{\leq k}^-:=\bigcup_{i=1}^{\min\{k,t\}}\mathcal K_i^-$ and let $\mathcal R_{\leq k}$ denote the set of all cycles in $\mathcal{K}_{\le k}^-$ which share some vertex with at least one of the cycles $\{R_1,\ldots,R_{\min\{k,t\}}\}$.

    We next plan to apply induction to $\mathcal K_{\leq k}^-$ and $\mathcal R_{\leq k}$ to obtain a new collection $\mathcal H_{\leq k}\subseteq \mathcal K_{\le k}^-$, as well as, for each $k+1\leq i\leq t$, to $\mathcal K_i^-$ and $\mathcal R_i$ to obtain a new collection $\mathcal H_i\subseteq \mathcal K_i^-$. 

    Note that by definition, we have $R_1 \in \mathcal C \setminus \bigcup_{i=1}^{t}\mathcal K_i^-$ and $R_i\in \mathcal C\setminus \mathcal K_i^-$ for every $i$, so $|\mathcal{K}_{\le k}^-|<|\mathcal C|$ and $|\mathcal{K}_i^-|<|\mathcal{C}|$ for every $k+1\le i \le t$. Hence, to ensure that these inductive applications are valid, it suffices to verify that for every $i\in [t]$, every connected component of $\mathcal{K}_i^-$ contains at least one cycle in $\mathcal R_i$. Indeed, since $\mathcal{R}_i\subseteq \mathcal{R}_{\le k}$ for every $i\in [\min\{k,t\}]$, this will then also imply that every connected component of $\mathcal{K}_{\le k}^-=\bigcup_{i=1}^{\min\{k,t\}}\mathcal{K}_i^-$ contains at least one cycle from $\mathcal{R}_{\le k}$. 

   So let $i\in [t]$ be given to us arbitrarily. Observe that $\mathcal K_i=\mathcal K_i^-\cup \{R_i\}$ is connected by definition. Hence, every connected component of $\mathcal K_i^-$ intersects $V(R_i)$ and, thus, contains a cycle of $\mathcal R_i$, as desired. 

   Thus, the above-mentioned inductive applications are justified and the arising collections $\mathcal{H}_{\le k}\subseteq \mathcal{K}_{\le k}^-$ and $\mathcal{H}_i\subseteq \mathcal{K}_i^-$ for $k+1\le i\le t$ are well-defined. Finally, we set\begin{align*}
        \mathcal H := \{R_1,\ldots,R_{\min\{k,t\}}\}\cup \mathcal H_{\leq k}\cup \bigcup_{i=k+1}^t\mathcal H_i.
    \end{align*}
     and claim that $\mathcal H$ satisfies the conditions \ref{cond:C1}--\ref{cond:C4} in the lemma. Note that $\mathcal{H}\cap \mathcal{R}=\{R_1,\ldots,R_{\min\{k,t\}}\}$ and so \ref{cond:C1} holds. Also, for every $i\in [\min\{k,t\}]$ we have that $\{R_i\}$ is vertex-disjoint from all cycles in $\mathcal H_{\leq k}\setminus \mathcal R_{\leq k}$ by the definition of $\mathcal R_{\leq k}$ and from $\bigcup_{i=k+1}^t\mathcal H_i$ by Claim~\ref{clm: disjointness}. But $\mathcal H_{\leq k}\cap \mathcal R_{\leq k}$ contains at most $k$ cycles by \ref{cond:C1}, so each $R_i$ intersects at most $\min\{k,t\}+k\le 2k$ cycles in $\mathcal{H}$. This verifies that \ref{cond:C2} holds for $\mathcal{H}$. We move on to \ref{cond:C3}. Consider any $C\in\mathcal H$, we want to prove that $C$ intersects at most $3k$ other cycles in $\mathcal{H}$. To do so, let us distinguish four cases. If $C\in\{R_1,\ldots,R_{\min\{k,t\}}\}$, then the statement follows from \ref{cond:C2}, which we just proved. If $C\in \mathcal{H}_{\le k}\cap \mathcal R_{\leq k}$, then $C$ intersects at most $2k$ cycles in $\mathcal H_{\leq k}$ by induction and $C$ does not intersect $\bigcup_{i=k+1}^t\mathcal H_i$ by Claim~\ref{clm: disjointness}. It follows that $C$ intersects at most $2k+\min\{k,t\}\le 3k$ cycles in $\mathcal H$. Next, suppose $C\in \mathcal H_{\leq k}\setminus\mathcal R_{\leq k}$. Then, $C$ intersects neither $\{R_1,\ldots,R_{\min\{k,t\}}\}$ (by definition of $\mathcal R_{\leq k}$) nor $\bigcup_{i=k+1}^t\mathcal H_i$ (by Claim~\ref{clm: disjointness}). Therefore, we are done by induction. Finally, suppose $C\in \mathcal H_i$ for some $k+1\leq i\leq t$. By Claim~\ref{clm: disjointness}, all cycles in $\mathcal{H}$ which $C$ intersects in some vertex must lie in $\mathcal H_i$, and by induction (using condition \ref{cond:C3}), there are at most $3k$ cycles in $\mathcal{H}_i$ which $C$ intersects. This completes the proof that $\mathcal{H}$ satisfies condition~\ref{cond:C3}.

    It is left to show that $\mathcal{H}$ also satisfies \ref{cond:C4}. 
    Recall that $V\subseteq V(\mathcal C)$ in the lemma statement was defined as the set of all vertices contained in connected components of $\mathcal C_p\cup \mathcal H$ which contain at least one cycle of $\mathcal{R}$. We then have to show that $\mathbb{E}[w(V)]\ge w(\mathcal C)^{1-\delta}$. 

    Using the definitions of $\mathcal{H},\mathcal{H}_{\le k},\mathcal{R}_{\le k}$ and $\mathcal{H}_i, \mathcal{R}_i$, we can observe that each of the following sets of vertices is contained in $V$:
    \begin{itemize}
        \item The connected components of $(\mathcal{K}_{\le k}^-)_p\cup \mathcal{H}_{\le k}$ which contain at least one cycle in $\mathcal{R}_{\le k}$ (note that these components intersect vertices of at least one $R_i$, $i \in [\min\{k,t\}]$, all of which we include in $\mathcal H$) form part of $V$.
        \item For each integer $i$ with $k+1\le i\le t$ and such that $R_i\in \mathcal{C}_p$ (which happens with probability $p$), we have that all connected components of $(\mathcal{K}_i^-)_p\cup \mathcal{H}_i$ which contain at least one cycle in $\mathcal{R}_i$ (and hence intersect $R_i\in \mathcal{R}$) are contained in $V$.
        \item We have $\bigcup_{i=1}^{\min\{k,t\}}V(R_i)\subseteq V$, and each vertex in $\bigcup_{i=k+1}^{t}V(R_i)$ is contained in $V$ with probability at least $p$. 
    \end{itemize}

Pause to note that by Claim~\ref{clm: disjointness}, the sets described in the first two bullets above are pairwise disjoint from one another. Furthermore, by inductive applications of~\ref{cond:C4}, we have that the expected total weight of vertices of the first type is at least $w(\mathcal{K}_{\le k}^-)^{1-\delta}$ and the expected total weight of the vertices of the second type is at least $p\sum_{i=k+1}^{t} w(\mathcal{K}_i^-)^{1-\delta}$. Hence, the expected total weight of vertices of the first and second types is at least
$$w(\mathcal{K}_{\le k}^-)^{1-\delta}+p\sum_{i=k+1}^{t} w(\mathcal{K}_i^-)^{1-\delta}.$$

Finally, observe that the additional weight provided by the third type of vertices described above is, in expectation, at least $w(\mathcal{K}_{\le k})-w(\mathcal{K}_{\le k}^-)+p\cdot w(V(\mathcal C)\setminus V\left(\mathcal{K}_{\le k}\cup \bigcup_{i=k+1}^t\mathcal{K}_i^-\right))$. 
All in all, we thus find that
\begin{align*}
        \mathbb E[w(V)]&\geq w(\mathcal{K}_{\le k}^-)^{1-\delta}+w(\mathcal{K}_{\le k})-w(\mathcal{K}_{\le k}^-)+p\left(\sum_{i=k+1}^{t} w(\mathcal{K}_i^-)^{1-\delta}\right)+p\cdot w\left(V(\mathcal C)\setminus V\left(\mathcal{K}_{\le k}\cup \bigcup_{i=k+1}^t\mathcal{K}_i^-\right)\right) \\
        &\ge w(\mathcal{K}_{\le k})^{1-\delta}+p\left(\left(\sum_{i=k+1}^{t} w(\mathcal{K}_i^-)^{1-\delta}\right)+w\left(V(\mathcal C)\setminus V\left(\mathcal{K}_{\le k}\cup \bigcup_{i=k+1}^t\mathcal{K}_i^-\right)\right)\right), 
    \end{align*}
    where we used the easily verified fact that $a^{1-\delta}+b\ge (a+b)^{1-\delta}$ for all $a,b\in\mathbb{Z}_{\ge 0}$ and $\delta \in (0,1)$.

    In the following, for every $i\in [t]$ let us denote $Y_i:=~V(\mathcal{K}_i)\setminus \bigcup_{j<i}V(\mathcal{K}_j)$. Note that by our definition of $\mathcal{K}_1,\ldots,\mathcal{K}_t$, we then have $w(Y_1)\ge w(Y_2)\ge \cdots\ge w(Y_t)$. Note in addition that the sets $Y_1,\ldots,Y_t$ form a partition of $V(\mathcal{C})$. We can then rewrite $V(\mathcal{K}_{\le k})=Y_1\cup\cdots\cup Y_{\min\{k,t\}}$. Further note that $\bigcup_{i=k+1}^t Y_i=V(\mathcal C)\setminus V(\mathcal{K}_{\le k})$ , and hence  $$V(\mathcal C)\setminus V\left(\mathcal{K}_{\le k}\cup \bigcup_{i=k+1}^t\mathcal{K}_i^-\right)=\bigcup_{i=k+1}^t Y_i\setminus \bigcup_{i=k+1}^{t}V(\mathcal{K}_i^-)=\bigcup_{i=k+1}^{t}(Y_i\setminus V(\mathcal{K}_i^-)),$$ where in the last step we used that $Y_i\cap V(\mathcal{K}_j^-)=\emptyset$ whenever $i\neq j\in [t]$, which follows from Claim~\ref{clm: disjointness} and the definition of $Y_1,\ldots,Y_t$. This allows us to deduce the following lower bound on $\mathbb{E}[w(V)]$:
\begin{align*}\mathbb{E}[w(V)]&\ge w(\mathcal{K}_{\le k})^{1-\delta}+p\left(\left(\sum_{i=k+1}^{t} w(\mathcal{K}_i^-)^{1-\delta}\right)+w\left(V(\mathcal C)\setminus V\left(\mathcal{K}_{\le k}\cup \bigcup_{i=k+1}^t\mathcal{K}_i^-\right)\right)\right)\\
    &=w(\mathcal{K}_{\le k})^{1-\delta}+p\left(\sum_{i=k+1}^{t} w(\mathcal{K}_i^-)^{1-\delta}+\sum_{i=k+1}^t w(Y_i\setminus V(\mathcal{K}_i^-))\right)\\
    &\ge w(\mathcal{K}_{\le k})^{1-\delta}+p\sum_{i=k+1}^t (w(\mathcal{K}_i^-)^{1-\delta}+w(Y_i\setminus V(\mathcal{K}_i^-)))\\
    &\ge w(\mathcal{K}_{\le k})^{1-\delta}+p\sum_{i=k+1}^t w(Y_i)^{1-\delta},
    \end{align*} 
    where we again used the inequality $a^{1-\delta}+b\ge (a+b)^{1-\delta}$ in the last step.

    
    Before further estimating the above expression, it will be useful to prove the following simple inequality for non-negative real numbers.

    \begin{fact}\label{fact: inequality}
    Let $q\in\mathbb{Z}_{\ge 0}$, $k \in \mathbb{N}$, $\delta\in(0,1)$, and $p\ge k^{-\delta}$ as well as $x, y_1,\ldots,y_q\geq 0$ be real numbers such that $y_i\le \frac{x}{k}$ for each $i\in [q]$. Then,
    $$
    x^{1-\delta}+p\sum_{i=1}^q y_i^{1-\delta}\geq \left(x+\sum_{i=1}^qy_i\right)^{1-\delta}.
    $$
\end{fact}
\begin{claimproof}
  The statement trivially holds when $x=0$, so moving on suppose $x>0$. We then have 
    \begin{align*}
        &x^{1-\delta}+p\sum_{i=1}^q y_i^{1-\delta}\geq x^{1-\delta}+p\sum_{i=1}^q y_i\cdot (x/k)^{-\delta}\\ &\ge \left(x+\sum_{i=1}^{q}y_i\right)\cdot x^{-\delta}\ge \left(x+\sum_{i=1}^qy_i\right)^{1-\delta}.
    \end{align*}    
    \vspace{-0.8cm}

\textcolor{white}{hack}
\end{claimproof}

Let now $x:=w(\mathcal{K}_{\le k})=w(Y_1)+\cdots+w(Y_{\min\{k,t\}})$, $q:=t-\min\{k,t\}$, and $(y_1,\ldots,y_{q}):=\linebreak(w(Y_{k+1}),\ldots,w(Y_t))$. Since $w(Y_i)\le w(Y_{\min\{k,t\}})\le \frac{1}{\min\{k,t\}}x$ for every $k+1\le i \le t$, we can see that the numbers $q,x,y_1,\ldots,y_q$ satisfy the conditions of Fact~\ref{fact: inequality} (also if $k>t$, since then $q=0$ and the conditions are trivially satisfied). Hence, we obtain
 \begin{align*}
        \mathbb{E}[w(V)]&\ge w(\mathcal K_{\leq k})^{1-\delta}+p\sum_{i=k+1}^t w(Y_i)^{1-\delta}=x^{1-\delta}+p\sum_{j=1}^q y_j^{1-\delta}\\
        &\ge \left(x+\sum_{j=1}^q y_j\right)^{1-\delta}=\left(\sum_{i=1}^tw(Y_i)\right)^{1-\delta}=w(\mathcal{C})^{1-\delta},
    \end{align*}
    as desired. This shows that $\mathcal{H}$ also satisfies condition~\ref{cond:C4} of the inductive assertion, concluding the proof of the lemma.
\end{proof}

\end{document}
